\chapter{Blastomycosis}
\index{Blastomycosis}

\section*{Definition and Overview}
Blastomycosis is an uncommon fungal infection caused by \textit{Blastomyces dermatitidis}, a dimorphic fungus that lives in moist soil, decaying wood, and leaf litter. The infection usually begins in the lungs after spores are inhaled, and in a proportion of people it spreads through the bloodstream to the skin, bones, joints, prostate, and, rarely, the central nervous system. Blastomycosis is not contagious from person to person, and it is not transmitted from animals to humans by direct contact, although infected dogs can act as sentinels for disease in the surrounding environment.

\section*{Epidemiology}
Blastomycosis is endemic in parts of North America, particularly the region around the Great Lakes, the St Lawrence River valley, and the Ohio, Mississippi, and Missouri river valleys, with additional foci in parts of Canada and Africa. Most people who become infected live in or have travelled through these areas, and sporadic outbreaks have been linked to construction, excavation, camping, and other activities that disturb soil and rotting wood. Men are affected more often than women, usually in middle age, and the infection is more common in people whose work or recreation brings frequent soil contact. Disease in people with impaired immunity tends to be more severe and more likely to involve multiple organs.

\section*{Aetiology and Pathophysiology}
\textit{Blastomyces dermatitidis} grows as a mould in the environment, producing spores that become airborne when soil or decaying matter is disturbed. When inhaled, the spores reach the alveoli, where the warmth of the body triggers conversion into the yeast form, which multiplies at body temperature.

\begin{itemize}
    \item \textbf{Inhalation of spores} from disturbed soil, rotting wood, or watercourses is the usual route of infection
    \item \textbf{Rural exposure} on farms, riverbanks, and woodland increases risk
    \item \textbf{Direct inoculation} of the skin through a contaminated wound occurs rarely
    \item \textbf{Host immunity is decisive:} cellular immunity controls the infection, so severe and disseminated disease is more common in people with weakened immune systems
\end{itemize}

In the lungs, the yeast provokes a granulomatous inflammatory response. Infection may resolve spontaneously in healthy people, cause acute or chronic pneumonia, or disseminate haematogenously to extrapulmonary sites. The skin, bones, joints, and genitourinary tract are the most commonly affected sites outside the lungs.

\section*{Clinical Features}
Symptoms usually appear two weeks to several months after exposure, and many infected people remain asymptomatic or experience only a mild, flu-like illness. Clinical patterns include acute pneumonia, chronic pulmonary infection that mimics tuberculosis or cancer, and disseminated disease.

\begin{itemize}
    \item \textbf{Respiratory:} fever, chills, night sweats, cough that may produce sputum or blood, pleuritic chest pain, and shortness of breath
    \item \textbf{General:} fatigue, muscle aches, anorexia, and weight loss
    \item \textbf{Skin:} verrucous, warty plaques, crusted ulcers, abscesses, or pustular lesions that heal poorly; often the first clue to the diagnosis
    \item \textbf{Musculoskeletal:} bone pain, joint swelling, and tenderness from osteomyelitis or arthritis
    \item \textbf{Genitourinary:} in men, prostatitis with urinary symptoms
    \item \textbf{Nervous system:} headache, confusion, or focal neurological signs in the rare meningeal or cerebral form
\end{itemize}

\section*{Differential Diagnosis}
The pulmonary and cutaneous features overlap with several common and serious conditions.

\begin{itemize}
    \item Community-acquired or aspiration pneumonia, and tuberculosis
    \item Influenza and other viral respiratory infections
    \item Bronchogenic carcinoma or metastatic lung disease
    \item Other endemic mycoses, especially histoplasmosis and coccidioidomycosis
    \item Chronic skin infections, pyoderma gangrenosum, and squamous cell carcinoma of the skin
    \item Sarcoidosis, in its chronic pulmonary and skin forms
\end{itemize}

\section*{When to Seek Medical Care}
Seek medical assessment for a persistent cough with fever, night sweats, or weight loss, skin lesions that do not heal, or bone and joint pain, particularly after travel to or residence in an endemic area.

\textbf{Contact your local emergency services for:}
\begin{itemize}
    \item Difficulty breathing or shortness of breath at rest
    \item Coughing up significant amounts of blood
    \item High fever with confusion, drowsiness, or a severe headache
    \item Signs of sepsis: rapid breathing, rapid heart rate, low blood pressure, or mottled skin
\end{itemize}

\section*{Diagnosis and Investigations}
A high index of suspicion is required, because blastomycosis is easily mistaken for more common infections. The diagnosis is confirmed by identifying the organism in tissue or culture.

\begin{itemize}
    \item \textbf{History:} residence, travel, occupation, and recreational exposure in endemic areas
    \item \textbf{Chest imaging:} X-ray or CT showing lobar or nodular pneumonia, cavitation, or mass-like lesions
    \item \textbf{Sputum microscopy and culture:} wet-mount examination may reveal the characteristic broad-based budding yeast, and culture confirms the species
    \item \textbf{Tissue biopsy:} of skin, lung, or bone lesions, with histopathology showing yeast forms
    \item \textbf{Antigen testing:} urinary and serum \textit{Blastomyces} antigen assays, which support the diagnosis but cross-react with other endemic mycoses
    \item \textbf{Lumbar puncture:} for cerebrospinal fluid examination when central nervous system involvement is suspected
    \item \textbf{Specialist referral:} to an infectious diseases physician for severe, disseminated, or treatment-resistant disease
\end{itemize}

\section*{Management}
Antifungal therapy is the mainstay, and early treatment of disseminated disease markedly improves outcomes.

\begin{itemize}
    \item \textbf{Itraconazole:} the oral drug of choice for mild to moderate pulmonary and non-central-nervous-system disease, usually continued for six to twelve months with monitoring of drug levels
    \item \textbf{Amphotericin B:} given intravenously for severe pneumonia, disease in immunocompromised patients, and central nervous system involvement, before stepping down to an oral triazole
    \item \textbf{Voriconazole or posaconazole:} alternatives for resistant or relapsing disease
    \item \textbf{Monitoring:} regular clinical and radiological review until symptoms resolve and infection is cleared
    \item \textbf{Supportive care:} rest, hydration, and antipyretics, together with drainage of abscesses where needed
\end{itemize}

\section*{Complications}
\begin{itemize}
    \item Severe pneumonia complicated by the acute respiratory distress syndrome
    \item Dissemination to skin, bones, joints, and the genitourinary tract
    \item Meningitis and cerebral abscess formation
    \item Abscesses and sinus formation in the lungs, skin, and bones
    \item Slow response or relapse after apparently successful treatment
    \item Progressive and life-threatening disease in people with weakened immunity
\end{itemize}

\section*{Prognosis and Outcomes}
With appropriate antifungal treatment, most otherwise healthy people recover fully, although therapy may need to continue for many months and some residual scarring of the lungs is common. Untreated disease tends to progress and can rarely be fatal, particularly in immunocompromised patients. Relapse is more likely in those with impaired immunity, in central nervous system disease, and when therapy is shortened or interrupted; these groups require particularly close follow-up.

\section*{Prevention and Self-Care}
Because the fungus is a natural inhabitant of soil, complete avoidance of exposure in endemic areas is impractical, but risk can be reduced.

\begin{itemize}
    \item Be aware of the risk when camping, farming, hunting, or working with soil and decaying wood in endemic regions
    \item Minimise disturbance of soil and waterlogged wood where possible
    \item Wear appropriate respiratory protection, such as an N95-style mask, when working in dusty soil-rich environments
    \item Wash skin wounds promptly and cover them when working outdoors
    \item Seek early medical review for a persistent cough or unusual skin lesions after exposure
    \item People with weakened immunity should avoid high-risk outdoor activities in endemic areas
\end{itemize}

\section*{Sources and Further Reading}
The following sources provide reliable, up-to-date information for patients and students of medicine.

\begin{itemize}
    \item Centers for Disease Control and Prevention. \textit{Blastomycosis -- causes, symptoms, and prevention.} https://www.cdc.gov/fungal/diseases/blastomycosis/
    \item Mayo Clinic. \textit{Overview of blastomycosis and its treatment.} https://www.mayoclinic.org/diseases-conditions/blastomycosis/
    \item MSD Manual. \textit{Blastomycosis -- professional edition.} https://www.msdmanuals.com/
    \item MedlinePlus. \textit{Blastomycosis -- patient information.} https://medlineplus.gov/
    \item World Health Organization. \textit{Fungal priority pathogens list.} https://www.who.int/
\end{itemize}
